\section{Exact positive phase fibers and their denominator boundary}
\label{pfs-section}

For a positive integer $B$, put
\[
 \mathcal A_B=\{3k:k\in\mathbb Z,\ B\leq k<2B\}.
\]
Since $(a+\zeta)(b+\zeta)=(ab-1)+(a+b+1)\zeta$, the phase of
an ordered pair is $(ab-1)/(a+b+1)$. An unordered pair below permits
$a=b$, counted once. These are actual root parameters; no common
deep-prime condition is implicit in the definition.

\begin{theorem}[Exact divisor parameterization]\label{pfs-divisors}
Let $u,v$ be coprime positive integers and $K=u^2+uv+v^2$.
Positive integer pairs of phase $u/v$ are in bijection with
positive divisors $P\mid K$ satisfying $P\equiv-u\pmod v$,
by
\begin{equation}\label{pfs-inverse}
 P=va-u,\quad Q=vb-u,\qquad
 a=\frac{P+u}{v},\quad b=\frac{K/P+u}{v}.
\end{equation}
For pairs with both coordinates divisible by three, the phase
necessarily satisfies
\begin{equation}\label{pfs-three}
 u+v\equiv0\pmod3,\qquad 3\nmid uv.
\end{equation}
Under \eqref{pfs-three}, the exact divisor congruence is
$P\equiv-u\pmod{3v}$. The additional conditions for the actual
block are
\begin{equation}\label{pfs-interval}
 3vB-u\leq P<6vB-u,\qquad
 3vB-u\leq K/P<6vB-u.
\end{equation}
\end{theorem}
\begin{proof}
The phase equation gives
\[
 (va-u)(vb-u)=K,\qquad
 b(va-u)=u(a+1)+v>0,\quad a(vb-u)=u(b+1)+v>0.
\]
Thus both factors are positive. Conversely the congruence for $P$
implies $\gcd(P,v)=1$. Since $K\equiv u^2\pmod v$, its complementary
divisor $Q$ automatically satisfies $Q\equiv-u\pmod v$.
Both reconstructed coordinates are positive integers, and expansion
of $PQ=K$ proves their phase equation. These constructions are inverse.

For coordinates divisible by three the unreduced numerator is $-1$
modulo three and its denominator is $1$ modulo three. Their gcd is
prime to three, so reduction proves \eqref{pfs-three}. Under that
condition $\gcd(u,3v)=1$ and $K-u^2=v(u+v)$ is divisible by $3v$.
The same complementary-divisor argument now works modulo $3v$.
The coordinate interval is exactly \eqref{pfs-interval}.
Before knowing that a candidate phase has a nonempty block fiber,
positivity of $P,Q$ must still be imposed. The next theorem proves
that its lower endpoint is positive whenever the fiber is nonempty.
\end{proof}

\begin{theorem}[A denominator-sensitive fiber bound]\label{pfs-denominator}
A nonempty block fiber of reduced phase $u/v$ satisfies
\begin{gather}
 \frac{9B^2-1}{6B+1}\leq\frac uv
 \leq\frac{(6B-3)^2-1}{12B-5},\qquad
 B<\frac uv<3B.\label{pfs-phase-range}
\end{gather}
Its unordered and ordered cardinalities satisfy respectively
\begin{align}
 N_{\rm unordered}&\leq1+\left\lfloor\frac{2B-2}{v}\right\rfloor,
 \label{pfs-unordered}\\
 N_{\rm ordered}&\leq
 \min\left\{\tau(K),\,2\left(1+\left\lfloor\frac{2B-2}{v}\right\rfloor\right)\right\}.
 \label{pfs-ordered}
\end{align}
In particular $v>2B-2$ implies at most one unordered pair.
\end{theorem}
\begin{proof}
At fixed $b$, the difference of phases at $a$ and $a'$ is
\[
 \frac{(a-a')(b^2+b+1)}{(a+b+1)(a'+b+1)}.
\]
Monotonicity in both coordinates proves the first two bounds of
\eqref{pfs-phase-range}. Its lower endpoint exceeds $B$ because
$3B^2-B-1>0$ for $B\geq1$. Also
$(ab-1)/(a+b+1)<ab/(a+b)<3B$ when $a,b<6B$.
Consequently $L=3vB-u>0$. With $U=6vB-u$, the interval condition
can equivalently be written
$L\leq P<U$ and $K/U<P\leq K/L$, with the displayed strict endpoints.

Put $s=a+b+1$. The reduced phase forces $v\mid s$, while the
block gives
\[
 6B+1\leq s\leq12B-5,\qquad s\equiv1\pmod3.
\]
Since $3\nmid v$, the possible sums form one progression of step
$3v$. The interval length is $6B-6$, so there are at most
$1+\lfloor(2B-2)/v\rfloor$ such sums. For each $s$, the product is
$ab=us/v+1$. The quadratic with this product and sum $s-1$
determines at most one unordered pair. This proves
\eqref{pfs-unordered}; ordering at most doubles the number, and
Theorem~\ref{pfs-divisors} supplies the bound $\tau(K)$.
Diagonal pairs contribute only one ordered pair. The proof includes
$B=1$ and both endpoints of the sum interval.
\end{proof}

This proves an actual subclass with an explicit denominator
condition. It does not establish that small-denominator fibers
are rare in a specified deep-root set.

\begin{proposition}[Two actual failures of stronger injectivity]
\label{pfs-examples}
For $n=7$, $B=n^4=2401$, phase $4526$ has the distinct unordered
pairs
\[
 (7809,10767),\quad(8397,9819).
\]
Phase $4799$ has at least three distinct unordered pairs,
\[
 (7668,12828),\quad(8922,10386),\quad(9480,9720).
\]
All coordinates belong to $\mathcal A_B$.
\end{proposition}
\begin{proof}
For the first phase,
\[
 K=20489203=7^2\cdot67\cdot79^2
   =3283\cdot6241=3871\cdot5293.
\]
The factors are $(a-4526,b-4526)$ in the two rows. For the second,
\begin{gather*}
 K=23035201=7\cdot19\cdot31\cdot37\cdot151,\\
 K=2869\cdot8029=4123\cdot5587=4681\cdot4921.
\end{gather*}
These factors are $(a-4799,b-4799)$. The inverse formulas in
Theorem~\ref{pfs-divisors} prove all phase equalities. Each
coordinate is a multiple of three in $[7203,14406)$.
\end{proof}

The first example rules out unconditional unordered injectivity;
the second rules out a universal bound of two unordered pairs.
Neither meets the large-denominator hypothesis of
Theorem~\ref{pfs-denominator}. They do not assert common
fourth-depth membership at a prime, and they do not refute the
signed-tail gate.

\begin{theorem}[A finite multiple-mark interface]\label{pfs-marks}
Let $S\subseteq\mathcal A_B$, and let $H$ be a finite set with
a binary operation. Assign images $\phi(a)\in H$ for $a\in S$.
Suppose an integer $m>864B^3$ has the property
\begin{equation}\label{pfs-div-input}
 \phi(a)\phi(b)=\phi(c)\phi(d)
 \quad\Longrightarrow\quad
 m\mid(ab-1)(c+d+1)-(cd-1)(a+b+1)
\end{equation}
for $a,b,c,d\in S$. For a positive integer $v_0$, let $E_{<v_0}$
count the ordered pairs in $S^2$ whose reduced phase denominator
is less than $v_0$. Then
\begin{equation}\label{pfs-low-denominator-remainder}
 |S|^2\leq E_{<v_0}
  +2|H|\left(1+\left\lfloor\frac{2B-2}{v_0}\right\rfloor\right).
\end{equation}
\end{theorem}
\begin{proof}
The actual degree-three determinant has absolute value at most
$864B^3$: each pair has coordinates between zero and $36B^2$,
and between zero and $12B$, respectively. Its divisibility by
$m$ therefore forces it to vanish. Every nonempty product-image
fiber has a single phase. Remove the low-denominator ordered pairs
and partition the rest by their image in $H$. Each occupied fiber
has denominator at least $v_0$, so Theorem~\ref{pfs-denominator}
gives its stated bound. Summing over $H$ proves the result.
No bound on $E_{<v_0}$ is presumed.
\end{proof}

For an application with several marked primes, let $q_i>3$ be
distinct, $q_i\nmid n$, and impose actual conditions
$q_i^{e_i}\mid T_n(a)$, with $e_i\geq1$, where $T_n(a)$ is
the absolute cubic boundary of $(a+\zeta)^n$. A hit forces
$N(a+\zeta)$ to be a unit modulo $q_i$: the cubic difference
would otherwise force both conjugates to vanish, while their
difference is a unit at $q_i>3$. Thus
$\alpha_a=(a+\zeta)/(a+\bar\zeta)$ is defined and belongs to
the $3n$-torsion subgroup $H_i$ modulo $q_i^{e_i}$.
The split or inert norm-one residue group is cyclic of order
$q_i-\chi_i$, where $\chi_i=1$ or $-1$, respectively. Simple
lifting of $X^{3n}-1$, using $q_i\nmid3n$, gives
\[
 |H_i|=\gcd(3n,q_i-\chi_i).
\]
Take $H=\prod_i H_i$ and $m=\prod_i q_i^{e_i}$. Equal pair images
give the conjugate-cross-product divisibility at every prime power,
and coprimality gives \eqref{pfs-div-input}. If $m>864B^3$,
Theorem~\ref{pfs-marks} applies.

This uses combined precision, but its target size is
$\prod_i|H_i|$, not one copy of $3n$. It yields no pointwise
far-tail bound by itself. Actual incidence of small phase
denominators, or additional dependence among product images,
remains an explicit question.

The accompanying exact integer replay checks $2870$ ordered pairs
from blocks $1\leq B\leq20$, their $1540$ distinct block fibers,
and $1599$ positive input pairs from $1\leq a,b\leq40$.
It fully factors the two displayed constants and enumerates their
divisors. A bounded search of $1201$ integer phases
$u\equiv2\pmod3$ in $[\lfloor3B/2\rfloor,3B)$ at $B=2401$
finds the displayed three-pair phase. This is not a search over all
rational phases and proves no uniform bound by testing.
The ordinary proofs above and the finite evidence are separate;
no Lean proof of the new denominator refinement is claimed.
